% Generated draft: requires mathematical and source review.
\documentclass[11pt]{article}
\usepackage[margin=25mm]{geometry}
\usepackage{fontspec}
\setmainfont{Latin Modern Roman}
\usepackage{amsmath,amssymb,mathtools,unicode-math}
\setmathfont{Latin Modern Math}
\usepackage{enumitem,needspace,xurl,hyperref}
\setlength{\parindent}{0pt}
\setlength{\parskip}{5pt}
\setlength{\emergencystretch}{3em}
\newcommand{\N}{\mathbb{N}}
\newcommand{\Z}{\mathbb{Z}}
\newcommand{\Q}{\mathbb{Q}}
\newcommand{\R}{\mathbb{R}}
\newcommand{\C}{\mathbb{C}}
\newcommand{\F}{\mathbb{F}}
\newcommand{\A}{\mathbb{A}}
\newcommand{\PP}{\mathbb{P}}
\newcommand{\E}{\mathbb{E}}
\newcommand{\eps}{\varepsilon}
\newcommand{\dd}{\,\mathrm{d}}
\begin{document}
\section*{UnsolvedMath ID: 11000026}
\textbf{Source problem code:} catalogue entry
\subsection{Definitions and mathematical statement}
Let $\operatorname{Teich}(\Sigma_g)$ be the Teichm\"uller space of a genus-$g$ surface with the source's marking and boundary conventions. Its Teichm\"uller distance is conventionally
\[
 d_{\mathrm T}(X,Y)=\tfrac12\inf_f\log K(f),
\]
where $f$ ranges over quasiconformal maps in the marking class and $K(f)$ is maximal dilatation. Question 3.5 asks whether this metric space is almost convex.

The supplied excerpt does not give the definition of almost convex for this continuous metric space. A Cannon-style metric formulation would fix a basepoint $o$ and scale $R>0$ and ask for $C_R<\infty$ such that, for every radius $r$ and all $x,y$ on the sphere about $o$ with $d_{\mathrm T}(x,y)\leq R$, there is a rectifiable path from $x$ to $y$ of length at most $C_R$ lying in the closed ball of radius $r$ about $o$. Whether the source requires this for one fixed scale, every scale, one basepoint, or uniformly over all basepoints is not supplied. This is therefore a conditional interpretation rather than an imposed definition.

For context the source attributes to Cannon [Ca] almost convexity of fundamental groups of closed negatively curved manifolds with respect to any generating set, understood in the usual finitely generated word-metric setting. It reports his conjectural extensions to finite-volume manifolds and to nonpositively curved manifolds. These discrete group statements are background and do not establish the Teichm\"uller metric assertion.
\subsection{Short English statement}
Is Teichm\"uller space almost convex in its Teichm\"uller metric? The exact bounded-path definition and its basepoint and scale quantifiers must be recovered from the source context.
\subsection{Sources}
\begin{itemize}
\item Local registry: \url{data/unsolvedmath/problems.json}, permanent UnsolvedMath ID $11000026$, exact source code \texttt{catalogue entry}.
\item Farb, \emph{Problems on Mapping Class Groups and Related Topics}, Question 3.5, PDF page 32: \url{https://www.math.uchicago.edu/~farb/papers/mcgbook.pdf}. Supplied access date: 2026-07-29; extraction: pdf-font-structure; page-disambiguated numbering among 329 headings.
\item The source cites Cannon [Ca] for almost-convex group results and proposed extensions. The bibliography entry and the preceding definition are not supplied. The registry reports no resolution located through 2026 without additional references.
\end{itemize}
\end{document}
